\documentclass[11pt]{article}
\usepackage[margin=1in]{geometry}
\usepackage{xcolor}
\usepackage[breakable,listings]{tcolorbox}
\tcbuselibrary{listings}
\usepackage{hyperref}
\usepackage{enumitem}
\usepackage{listings}
\usepackage{amsmath}
\usepackage{booktabs}
\usepackage{array}

% Define colors
\definecolor{osaorange}{RGB}{220,115,38}
\definecolor{aiblue}{RGB}{30,144,255}

% Configure hyperref
\hypersetup{
    colorlinks=true,
    linkcolor=blue,
    urlcolor=blue
}

% Configure R code listings
\lstset{
    language=R,
    basicstyle=\ttfamily\small,
    breaklines=true,
    breakatwhitespace=true,
    postbreak=\mbox{\textcolor{red}{$\hookrightarrow$}\space},
    frame=single,
    backgroundcolor=\color{gray!10},
    numbers=left,
    numberstyle=\tiny,
    keywordstyle=\color{blue},
    commentstyle=\color{green!60!black},
    stringstyle=\color{red},
    columns=flexible,
    keepspaces=true,
    xleftmargin=0pt,
    xrightmargin=0pt,
    linewidth=\textwidth,
    aboveskip=\baselineskip,
    belowskip=\baselineskip,
    framexleftmargin=0pt,
    framexrightmargin=0pt
}

\title{}
\author{}
\date{}

\begin{document}

% Header box
\begin{tcolorbox}[colback=osaorange,colframe=black,boxrule=2pt,arc=0pt,left=3pt,right=3pt,top=6pt,bottom=6pt,boxsep=2pt]
\centering
\textcolor{white}{\textbf{\Large H524 Introduction to Biostatistics}}\\
\textcolor{white}{\textbf{\Large Project Preparation Homework}}\\
\textcolor{white}{\textbf{\Large GROUP ASSIGNMENT}}\\
\textcolor{white}{\textbf{\Large Fall 2025}}
\end{tcolorbox}

\vspace{8pt}

\noindent\textbf{Due Date:} End of Week 8\\
\textbf{Total Points:} 100\\
\textbf{Format:} GROUP SUBMISSION - One PDF per group via Canvas\\
\textbf{Group Size:} 4-5 students

\section*{Overview}

This assignment marks the beginning of your group final project for H524. By now, you have learned descriptive statistics, probability, confidence intervals, hypothesis testing (t-tests, ANOVA), and nonparametric methods. This assignment asks you to:

\begin{enumerate}
\item Form a research team of 4-5 students
\item Select a biostatistics dataset of interest
\item Develop testable research questions
\item Conduct preliminary data exploration
\item Create an analysis plan using methods from weeks 1-8
\end{enumerate}

This preparation will serve as the foundation for your final project presentation and report at the end of the term.

\section*{Learning Objectives}

By completing this assignment, you will:
\begin{itemize}
\item Practice collaborative research planning and teamwork
\item Apply course concepts to real biomedical/public health data
\item Develop testable hypotheses appropriate for your dataset
\item Conduct comprehensive exploratory data analysis
\item Create reproducible, well-documented R code
\item Communicate statistical findings clearly
\end{itemize}

\section*{Timeline and Checkpoints}

\begin{table}[h]
\centering
\begin{tabular}{lp{10cm}}
\toprule
\textbf{Week} & \textbf{Milestone} \\
\midrule
Week 7 & Form groups; begin dataset exploration \\
Week 7.5 & Optional: Dataset approval office hours \\
Week 8 & Complete data exploration and analysis plan \\
Week 8 (end) & \textbf{SUBMISSION DEADLINE} \\
\midrule
Weeks 9-10 & Continue analysis (after learning new methods) \\
Week 10 (end) & Final project presentations and reports due \\
\bottomrule
\end{tabular}
\end{table}

\newpage

\section{Part A: Group Formation and Organization (15 points)}

\subsection{Required Information}

Submit the following information about your research team:

\subsubsection{Team Members (5 points)}

Complete the following table:

\begin{center}
\begin{tabular}{|l|l|l|}
\hline
\textbf{Name} & \textbf{Email} & \textbf{Primary Role(s)} \\
\hline
& & \\[2ex]
\hline
& & \\[2ex]
\hline
& & \\[2ex]
\hline
& & \\[2ex]
\hline
& & \\[2ex]
\hline
\end{tabular}
\end{center}

\textbf{Roles} (students may hold multiple roles):
\begin{itemize}
\item \textbf{Project Coordinator}: Schedules meetings, tracks deadlines, ensures all deliverables are complete
\item \textbf{Data Analyst}: Leads R coding, data cleaning, statistical analysis
\item \textbf{Writer/Editor}: Coordinates writing of reports, ensures clarity and grammar
\item \textbf{Methods Specialist}: Ensures appropriate statistical methods are used correctly
\item \textbf{Presenter}: Leads development of final presentation materials
\end{itemize}

\subsubsection{Communication and Collaboration Plan (5 points)}

Describe your group's plan for working together:

\begin{enumerate}[label=\alph*)]
\item \textbf{Primary communication method:} (e.g., email, Slack, Discord, text)
\item \textbf{Meeting schedule:} How often will you meet? When?
\item \textbf{File sharing:} How will you share code and documents? (e.g., Google Drive, GitHub, Canvas)
\item \textbf{Conflict resolution:} How will your group handle disagreements or unequal contributions?
\end{enumerate}

\subsubsection{Individual Contribution Statement (5 points)}

At the end of your submission, include a brief statement from EACH group member describing:
\begin{itemize}
\item What you personally contributed to this assignment
\item Approximately what percentage of the work you completed
\item Any challenges your group faced in collaboration
\end{itemize}

\textbf{Note:} All group members should contribute meaningfully. If there are significant disparities in effort, please discuss with the instructor.

\newpage

\section{Part B: Dataset Selection and Description (20 points)}

\subsection{Dataset Selection}

Your group must select ONE dataset for analysis. You have two options:

\subsubsection{Option 1: Choose from Curated Datasets (Recommended)}

The instructor has prepared several high-quality biostatistics datasets appropriate for this course:

\begin{enumerate}
\item \textbf{Framingham Heart Study (subset)}: Cardiovascular disease risk factors
\item \textbf{NHANES (subset)}: National health and nutrition data
\item \textbf{COVID-19 Clinical Outcomes}: Patient characteristics and outcomes
\item \textbf{Birth Weight Study}: Factors affecting infant birth weight
\item \textbf{Cancer Survival Data}: Treatment outcomes and prognostic factors
\item \textbf{Mental Health Survey}: Depression, anxiety, and demographic factors
\end{enumerate}

\textit{Detailed dataset descriptions and download links will be provided on Canvas.}

\subsubsection{Option 2: Propose Your Own Dataset (Requires Approval)}

If your group wants to use a different dataset, you must:

\begin{itemize}
\item Submit a 1-page justification by \textbf{Week 7, Wednesday}
\item Include: dataset source, variables available, number of observations
\item Explain why this dataset is suitable for H524 methods
\item Receive instructor approval before proceeding
\end{itemize}

\textbf{Dataset Requirements (ANY dataset must meet these criteria):}
\begin{itemize}
\item At least 100 observations
\item At least 8-10 variables (mix of continuous and categorical)
\item Publicly available OR you have permission to use it
\item Relevant to biostatistics, public health, or medicine
\item Suitable for methods covered in weeks 1-8 (and weeks 9-10)
\item Complete data documentation available
\end{itemize}

\subsection{Dataset Description (20 points)}

Provide a comprehensive description of your chosen dataset:

\subsubsection{Source and Context (5 points)}

\begin{itemize}
\item Full citation for the dataset (APA format)
\item Brief background: What was the original study purpose?
\item Population studied (who, where, when)
\item Study design (cross-sectional, cohort, experimental, etc.)
\end{itemize}

\subsubsection{Dataset Overview (10 points)}

\begin{itemize}
\item Total number of observations (rows)
\item Total number of variables (columns)
\item Brief description of each variable including:
  \begin{itemize}
  \item Variable name
  \item Variable type (continuous, discrete, categorical nominal, categorical ordinal)
  \item Units of measurement (if applicable)
  \item Possible range or categories
  \end{itemize}
\end{itemize}

\textbf{Create a table summarizing all variables}. Example format:

\begin{center}
\begin{tabular}{|l|l|p{4cm}|l|}
\hline
\textbf{Variable} & \textbf{Type} & \textbf{Description} & \textbf{Range/Categories} \\
\hline
age & Continuous & Patient age at enrollment & 18-85 years \\
\hline
sex & Nominal & Biological sex & Male, Female \\
\hline
bp\_category & Ordinal & Blood pressure category & Normal, Elevated, High \\
\hline
\end{tabular}
\end{center}

\subsubsection{Missing Data Assessment (5 points)}

\begin{itemize}
\item Which variables have missing data?
\item What percentage of observations are missing for each variable?
\item How will you handle missing data in your analysis?
\end{itemize}

\subsubsection{Why This Dataset? (Bonus question, not graded)}

In 2-3 sentences, explain what interested your group about this dataset and why it's important for public health or biomedicine.

\newpage

\section{Part C: Research Questions and Hypotheses (25 points)}

Develop 2-3 specific, testable research questions based on your dataset. Each research question should include:

\subsection{Research Question Format}

For EACH research question (2-3 total), provide:

\subsubsection{Question Statement (5 points per question)}

State your research question clearly. Good research questions are:
\begin{itemize}
\item \textbf{Specific}: Focus on particular variables
\item \textbf{Testable}: Can be answered with statistical methods
\item \textbf{Relevant}: Important for public health or clinical practice
\end{itemize}

\textbf{Examples of good research questions:}
\begin{itemize}
\item ``Is there a significant difference in mean systolic blood pressure between smokers and non-smokers?''
\item ``Are cholesterol levels normally distributed in this population?''
\item ``Is there an association between BMI category and diabetes diagnosis?''
\end{itemize}

\textbf{Examples of poor research questions:}
\begin{itemize}
\item ``What can we learn from this data?'' (too vague)
\item ``How does everything relate to everything?'' (not specific)
\item ``What causes heart disease?'' (too broad, requires complex modeling)
\end{itemize}

\subsubsection{Formal Hypotheses (5 points per question)}

For each research question, state:
\begin{itemize}
\item \textbf{Null hypothesis} ($H_0$): Statement of no effect/no difference
\item \textbf{Alternative hypothesis} ($H_A$ or $H_1$): What you expect to find
\end{itemize}

Use proper statistical notation where appropriate.

\textbf{Example:}
\begin{itemize}
\item $H_0$: $\mu_{\text{smokers}} = \mu_{\text{non-smokers}}$ (no difference in mean BP)
\item $H_A$: $\mu_{\text{smokers}} \neq \mu_{\text{non-smokers}}$ (two-sided alternative)
\end{itemize}

\subsubsection{Public Health Significance (5 points per question)}

For each research question, explain:
\begin{itemize}
\item Why is this question important?
\item What are the potential implications if your hypothesis is supported?
\item How might results inform public health policy or clinical practice?
\end{itemize}

Limit to 3-5 sentences per question.

\subsubsection{Planned Statistical Method (5 points per question)}

For each research question, specify:
\begin{itemize}
\item Which statistical test/method you will use (based on weeks 1-8 content)
\item Why this method is appropriate for your variables and question
\item What assumptions you need to check
\end{itemize}

\textbf{Methods you have learned (weeks 1-8):}
\begin{itemize}
\item Descriptive statistics and visualization
\item One-sample t-test
\item Two-sample t-test (independent samples)
\item Paired t-test
\item One-way ANOVA
\item Nonparametric tests (Mann-Whitney U, Kruskal-Wallis, Wilcoxon signed-rank)
\item Correlation (Pearson, Spearman)
\end{itemize}

\newpage

\section{Part D: Initial Data Exploration (25 points)}

Conduct a thorough exploratory data analysis (EDA) of your dataset. This section should include:

\subsection{Summary Statistics (10 points)}

For your key variables of interest:

\subsubsection{Continuous Variables}
Provide:
\begin{itemize}
\item Sample size (n)
\item Mean and median
\item Standard deviation and IQR
\item Minimum and maximum values
\item Any outliers identified
\end{itemize}

\textbf{Create summary tables}. Example:

\begin{center}
\begin{tabular}{lcccccc}
\toprule
\textbf{Variable} & \textbf{n} & \textbf{Mean} & \textbf{SD} & \textbf{Median} & \textbf{IQR} & \textbf{Range} \\
\midrule
Age (years) & 500 & 52.3 & 12.1 & 51.0 & 44-60 & 18-85 \\
SBP (mmHg) & 500 & 128.5 & 18.2 & 125.0 & 115-140 & 90-195 \\
\bottomrule
\end{tabular}
\end{center}

\subsubsection{Categorical Variables}
Provide:
\begin{itemize}
\item Frequency counts for each category
\item Percentages
\item Mode (most common category)
\end{itemize}

\textbf{Create frequency tables}. Example:

\begin{center}
\begin{tabular}{lcc}
\toprule
\textbf{Category} & \textbf{Count} & \textbf{Percentage} \\
\midrule
Male & 245 & 49.0\% \\
Female & 255 & 51.0\% \\
\textbf{Total} & \textbf{500} & \textbf{100.0\%} \\
\bottomrule
\end{tabular}
\end{center}

\subsection{Data Visualization (10 points)}

Create appropriate visualizations for your key variables. Include AT LEAST:

\begin{itemize}
\item \textbf{2-3 histograms} for continuous variables
  \begin{itemize}
  \item Assess distribution shape (normal, skewed, bimodal?)
  \item Include reference lines for mean/median if helpful
  \end{itemize}
\item \textbf{2-3 box plots} comparing groups
  \begin{itemize}
  \item Show distribution differences between categories
  \item Identify outliers visually
  \end{itemize}
\item \textbf{1-2 bar charts} for categorical variables
\item \textbf{1-2 scatter plots} if examining relationships between continuous variables
\end{itemize}

\textbf{All plots must:}
\begin{itemize}
\item Have clear titles
\item Have labeled axes with units
\item Be large enough to read easily
\item Be generated in R (no screenshots from other software)
\end{itemize}

\subsection{Distribution Assessment (5 points)}

For each continuous variable you plan to analyze:

\begin{itemize}
\item Assess normality visually (histogram, Q-Q plot)
\item Comment on skewness or outliers
\item Discuss whether parametric tests are appropriate or if you need nonparametric alternatives
\end{itemize}

\textbf{Example discussion:}

``Age appears approximately normally distributed based on the histogram, with no extreme outliers. The Q-Q plot shows points falling close to the reference line. Therefore, parametric methods (t-tests, ANOVA) are appropriate for this variable.

Systolic blood pressure shows slight right skewness with 3 potential outliers above 180 mmHg. We will verify these are not data entry errors. If skewness persists, we may use Mann-Whitney U test instead of two-sample t-test.''

\newpage

\section{Part E: Preliminary Analysis Plan (15 points)}

\subsection{Statistical Methods Selection (8 points)}

For each research question from Part C, describe in detail:

\begin{enumerate}
\item \textbf{Which test you will use} (e.g., two-sample t-test, one-way ANOVA, Mann-Whitney U)
\item \textbf{Why this test is appropriate}
  \begin{itemize}
  \item What are your independent and dependent variables?
  \item What type of data do you have?
  \item Does the test match your research question?
  \end{itemize}
\item \textbf{Assumptions to check}
  \begin{itemize}
  \item Independence of observations
  \item Normality (if parametric)
  \item Equal variances (if required)
  \item Sample size requirements
  \end{itemize}
\item \textbf{How you will check assumptions}
  \begin{itemize}
  \item Visual methods (histograms, Q-Q plots, box plots)
  \item Statistical tests (Shapiro-Wilk, Levene's test)
  \item What you will do if assumptions are violated
  \end{itemize}
\end{enumerate}

\subsection{Potential Challenges and Solutions (4 points)}

Identify potential issues you foresee and how you plan to address them:

\begin{itemize}
\item \textbf{Missing data}: How much is missing? Will you use complete case analysis, imputation, or another approach?
\item \textbf{Outliers}: How will you handle extreme values? Remove? Transform? Use robust methods?
\item \textbf{Small sample sizes}: Do you have adequate power for your planned tests?
\item \textbf{Assumption violations}: What alternatives will you use if parametric assumptions fail?
\item \textbf{Multiple comparisons}: Are you conducting multiple tests? Do you need to adjust p-values?
\end{itemize}

\subsection{Analysis Timeline (3 points)}

Create a brief timeline for completing your final project:

\begin{center}
\begin{tabular}{|l|p{9cm}|}
\hline
\textbf{Week} & \textbf{Tasks} \\
\hline
Week 8 & Complete data exploration (THIS ASSIGNMENT) \\
\hline
Week 9 & Learn contingency tables/chi-square; apply to dataset if relevant \\
\hline
Week 9 & Conduct hypothesis tests; create results tables \\
\hline
Week 10 & Learn regression; apply to dataset if relevant \\
\hline
Week 10 & Finalize analysis; create visualizations \\
\hline
Week 10 & Write final report; create presentation slides \\
\hline
Week 10 (end) & Practice presentation; submit all materials \\
\hline
\end{tabular}
\end{center}

\newpage

\section{Part F: R Code and Reproducibility (Appendix)}

\subsection{Code Organization}

Submit well-organized, fully commented R code that performs all analyses described in Parts B-E.

\subsubsection{Required Code Sections}

Your R code must include:

\begin{enumerate}
\item \textbf{Header Comments}
\begin{lstlisting}
# H524 Project Preparation - Group [Name]
# Date: [Date]
# Group Members: [Names]
# Dataset: [Dataset name]
# Purpose: Exploratory data analysis and research question development
\end{lstlisting}

\item \textbf{Package Loading}
\begin{lstlisting}
# Load required packages
library(tidyverse)  # For data manipulation and visualization
library(here)       # For file path management
# Add other packages as needed
\end{lstlisting}

\item \textbf{Data Import and Cleaning}
\begin{lstlisting}
# Import dataset
data <- read.csv("your_dataset.csv")

# Check structure
str(data)
summary(data)

# Handle missing data
# [Your code here]

# Check for duplicates
# [Your code here]

# Data type conversions if needed
# [Your code here]
\end{lstlisting}

\item \textbf{Exploratory Data Analysis}
\begin{lstlisting}
# Summary statistics for continuous variables
# [Your code here]

# Frequency tables for categorical variables
# [Your code here]

# Create visualizations
# Histograms
# Box plots
# Scatter plots
# [Your code here]
\end{lstlisting}

\item \textbf{Assumption Checking}
\begin{lstlisting}
# Normality assessment
# Q-Q plots
# Shapiro-Wilk test
# [Your code here]

# Check for outliers
# [Your code here]

# Variance homogeneity (if needed)
# [Your code here]
\end{lstlisting}

\item \textbf{Preliminary Statistical Tests} (optional at this stage)
\begin{lstlisting}
# If you want to run preliminary tests, include them here
# Document results
# [Your code here]
\end{lstlisting}
\end{enumerate}

\subsection{Code Quality Requirements}

Your code will be evaluated on:

\begin{itemize}
\item \textbf{Reproducibility (5 points)}
  \begin{itemize}
  \item Can the instructor run your code without errors?
  \item Are file paths handled properly? (Use relative paths or \texttt{here()} package)
  \item Are all required packages listed?
  \item Does the code produce all tables and figures in your report?
  \end{itemize}

\item \textbf{Organization and Clarity (5 points)}
  \begin{itemize}
  \item Is code well-organized into logical sections?
  \item Are variable names meaningful?
  \item Is there adequate white space?
  \end{itemize}

\item \textbf{Documentation (5 points)}
  \begin{itemize}
  \item Are there comments explaining what each section does?
  \item Are complex operations explained?
  \item Would another student understand your code?
  \end{itemize}
\end{itemize}

\subsection{RMarkdown Template (RECOMMENDED)}

We strongly recommend using RMarkdown to integrate your code and written report. A template is provided below:

\begin{lstlisting}
---
title: "H524 Project Preparation"
author: "Group Name: [Your names]"
date: "`r Sys.Date()`"
output:
  pdf_document:
    toc: true
    number_sections: true
---

# Part A: Group Information

[Your content here]

# Part B: Dataset Description

[Your content here]

```{r setup, include=FALSE}
knitr::opts_chunk$set(echo = TRUE, warning = FALSE, message = FALSE)
library(tidyverse)
```

```{r load-data}
# Load your dataset
data <- read.csv("your_dataset.csv")
```

# Part C: Research Questions

[Your content here]

# Part D: Data Exploration

## Summary Statistics

```{r summary-stats}
# Your summary statistics code
```

## Visualizations

```{r histograms}
# Create histograms
```

# Part E: Analysis Plan

[Your content here]

# Appendix: Full R Code

```{r ref.label=knitr::all_labels(), echo=TRUE, eval=FALSE}
# This automatically includes all code chunks
```
\end{lstlisting}

\newpage

\section*{Grading Rubric (100 points total)}

\begin{table}[h]
\centering
\begin{tabular}{lcl}
\toprule
\textbf{Component} & \textbf{Points} & \textbf{Key Criteria} \\
\midrule
\textbf{Part A: Group Organization} & 15 & \\
\quad Team information & 5 & Complete roster with clear roles \\
\quad Collaboration plan & 5 & Concrete meeting/communication plan \\
\quad Individual contributions & 5 & Each member provides honest statement \\
\midrule
\textbf{Part B: Dataset Description} & 20 & \\
\quad Source and context & 5 & Complete citation and background \\
\quad Variable documentation & 10 & Comprehensive table with all variables \\
\quad Missing data assessment & 5 & Identification and handling plan \\
\midrule
\textbf{Part C: Research Questions} & 25 & \\
\quad Question clarity & 5 & Specific, testable, relevant (per question) \\
\quad Hypotheses formulation & 5 & Proper null/alternative statements \\
\quad Public health significance & 5 & Clear justification of importance \\
\quad Method selection & 5 & Appropriate test chosen with rationale \\
\midrule
\textbf{Part D: Data Exploration} & 25 & \\
\quad Summary statistics & 10 & Complete tables for key variables \\
\quad Visualizations & 10 & Multiple appropriate plots with labels \\
\quad Distribution assessment & 5 & Thoughtful evaluation of assumptions \\
\midrule
\textbf{Part E: Analysis Plan} & 15 & \\
\quad Methods justification & 8 & Detailed plan with assumption checks \\
\quad Challenges identified & 4 & Realistic issues with solutions \\
\quad Timeline & 3 & Feasible schedule for completion \\
\midrule
\textbf{Part F: R Code Quality} & 15 & \\
\quad Reproducibility & 5 & Code runs without errors \\
\quad Organization & 5 & Clear structure, good naming \\
\quad Documentation & 5 & Adequate comments throughout \\
\midrule
\textbf{Professional Presentation} & -5 & Deductions for poor formatting, typos \\
\bottomrule
\textbf{TOTAL} & \textbf{100} & \\
\bottomrule
\end{tabular}
\end{table}

\subsection*{Additional Evaluation Notes}

\begin{itemize}
\item \textbf{Group work expectations}: All members should contribute meaningfully. Individual contribution statements will be reviewed.
\item \textbf{Statistical rigor}: Methods must be appropriate for your data and research questions.
\item \textbf{Clarity of writing}: Scientific writing should be clear, concise, and professional.
\item \textbf{Reproducibility}: Your code must run on the instructor's computer with only minor path changes.
\end{itemize}

\newpage

\section*{AI Usage Guidelines}

\begin{tcolorbox}[colback=aiblue!10,colframe=aiblue,title=\textbf{AI-Enhanced Learning for Group Projects}]

This is a GROUP assignment where collaboration and learning are paramount. AI tools can assist your learning, but you must understand and verify everything you submit.

\subsection*{Encouraged AI Uses:}

\begin{itemize}
\item[+] \textbf{Code debugging}: ``Why is this R code giving an error?''
\item[+] \textbf{Syntax help}: ``How do I create a grouped box plot in ggplot2?''
\item[+] \textbf{Concept clarification}: ``Explain the difference between parametric and nonparametric tests''
\item[+] \textbf{Writing improvement}: ``Help me make this explanation clearer'' (after you write a draft)
\item[+] \textbf{Method selection}: ``I have continuous outcome and 3 groups. What test should I use?''
\end{itemize}

\subsection*{Prohibited AI Uses:}

\begin{itemize}
\item[-] \textbf{Complete delegation}: Asking AI to write your entire analysis
\item[-] \textbf{Hypothesis generation}: AI should not choose your research questions
\item[-] \textbf{Interpretation without understanding}: Copying AI explanations you don't comprehend
\item[-] \textbf{Skipping group discussion}: Using AI instead of discussing with teammates
\end{itemize}

\subsection*{Required AI Documentation:}

If you use AI tools significantly, include a brief section at the end of your report:

\begin{itemize}
\item Which AI tool(s) did you use? (e.g., ChatGPT, Claude, GitHub Copilot)
\item For what purposes? (e.g., debugging code, explaining statistical concepts)
\item How did you verify the AI's suggestions?
\item What did you learn from the AI interaction?
\end{itemize}

\subsection*{Remember:}

\begin{itemize}
\item You are responsible for all content you submit
\item You must understand every line of code and every statistical interpretation
\item AI can assist learning, but cannot replace learning
\item The goal is to develop YOUR group's analytical skills
\end{itemize}

\end{tcolorbox}

\newpage

\section*{Submission Requirements}

\subsection*{What to Submit (One submission per group)}

\begin{enumerate}
\item \textbf{Main Report (PDF)}
  \begin{itemize}
  \item Generated from RMarkdown (strongly recommended) OR
  \item Word document converted to PDF OR
  \item LaTeX-generated PDF
  \item Must include all Parts A-E
  \item Must include all tables, figures, and narrative text
  \item Expected length: 8-12 pages (not including code appendix)
  \end{itemize}

\item \textbf{R Code File(s)}
  \begin{itemize}
  \item Submit .Rmd file (if using RMarkdown) OR
  \item Submit .R script file with all code
  \item Ensure code is reproducible
  \item Include comments throughout
  \end{itemize}

\item \textbf{Dataset File}
  \begin{itemize}
  \item If using a curated dataset: Include citation/link (no need to upload)
  \item If using custom dataset: Upload data file (CSV format preferred)
  \item Include data dictionary if not already documented
  \end{itemize}

\item \textbf{Individual Contribution Statements}
  \begin{itemize}
  \item Can be included at end of main PDF OR
  \item Submit as separate document
  \item Each group member must provide their own statement
  \end{itemize}
\end{enumerate}

\subsection*{File Naming Convention}

Use clear, consistent file names:

\begin{itemize}
\item Main report: \texttt{H524\_ProjectPrep\_GroupName.pdf}
\item R code: \texttt{H524\_ProjectPrep\_GroupName.R} or \texttt{.Rmd}
\item Dataset: \texttt{H524\_ProjectPrep\_GroupName\_data.csv}
\end{itemize}

\subsection*{Canvas Submission}

\begin{itemize}
\item ONE group member submits all files
\item All other group members should be listed in the Canvas submission
\item Submit by 11:59 PM on the due date (end of Week 8)
\item Late penalty: 10\% per day (unless prior arrangement made)
\end{itemize}

\newpage

\section*{Example Research Question Development}

To help you develop strong research questions, here are examples using a hypothetical cardiovascular disease dataset:

\subsection*{Example 1: Comparing Two Groups}

\textbf{Research Question:}
Is there a significant difference in mean systolic blood pressure between patients with and without diabetes?

\textbf{Hypotheses:}
\begin{itemize}
\item $H_0$: $\mu_{\text{diabetes}} = \mu_{\text{no diabetes}}$ (no difference in mean SBP)
\item $H_A$: $\mu_{\text{diabetes}} \neq \mu_{\text{no diabetes}}$ (there is a difference)
\end{itemize}

\textbf{Variables:}
\begin{itemize}
\item Independent: Diabetes status (categorical: yes/no)
\item Dependent: Systolic blood pressure (continuous, mmHg)
\end{itemize}

\textbf{Statistical Method:}
Two-sample t-test (or Mann-Whitney U if SBP not normally distributed)

\textbf{Assumptions to Check:}
\begin{itemize}
\item Independence: Patients are independent observations
\item Normality: Check histogram and Q-Q plot for SBP in each group
\item Equal variances: Levene's test or examine box plots
\end{itemize}

\textbf{Public Health Significance:}
Understanding blood pressure differences by diabetes status can inform screening recommendations and highlight the need for intensive BP management in diabetic patients.

\subsection*{Example 2: Comparing Multiple Groups}

\textbf{Research Question:}
Do mean cholesterol levels differ significantly across three BMI categories (normal weight, overweight, obese)?

\textbf{Hypotheses:}
\begin{itemize}
\item $H_0$: $\mu_{\text{normal}} = \mu_{\text{overweight}} = \mu_{\text{obese}}$ (all group means equal)
\item $H_A$: At least one group mean differs from the others
\end{itemize}

\textbf{Variables:}
\begin{itemize}
\item Independent: BMI category (categorical: 3 levels)
\item Dependent: Total cholesterol (continuous, mg/dL)
\end{itemize}

\textbf{Statistical Method:}
One-way ANOVA (or Kruskal-Wallis if assumptions violated)

\textbf{Assumptions to Check:}
\begin{itemize}
\item Independence: Patients are independent
\item Normality: Check histograms for cholesterol in each BMI group
\item Equal variances: Levene's test; examine box plots
\end{itemize}

\textbf{Public Health Significance:}
If higher BMI is associated with elevated cholesterol, this supports weight management interventions for cardiovascular disease prevention.

\subsection*{Example 3: Correlation/Association}

\textbf{Research Question:}
Is there a significant correlation between age and resting heart rate in this population?

\textbf{Hypotheses:}
\begin{itemize}
\item $H_0$: $\rho = 0$ (no correlation between age and heart rate)
\item $H_A$: $\rho \neq 0$ (there is a correlation)
\end{itemize}

\textbf{Variables:}
\begin{itemize}
\item Variable 1: Age (continuous, years)
\item Variable 2: Heart rate (continuous, beats per minute)
\end{itemize}

\textbf{Statistical Method:}
Pearson correlation (if both variables normally distributed) or Spearman correlation (if not)

\textbf{Assumptions to Check:}
\begin{itemize}
\item Linearity: Create scatter plot
\item Normality: Check histograms for both variables
\item Outliers: Identify extreme values that might influence correlation
\end{itemize}

\textbf{Public Health Significance:}
Understanding age-related changes in heart rate helps establish normal ranges and identify patients with abnormal cardiovascular responses.

\subsection*{Tips for Your Group}

\begin{itemize}
\item Start with simple, clear questions
\item Make sure your questions can be answered with methods from weeks 1-8
\item Choose questions that genuinely interest your group
\item Consider clinical or public health relevance
\item Ensure your dataset has the necessary variables
\end{itemize}

\newpage

\section*{Required R Packages}

Your analysis will likely use the following R packages. Install them before beginning:

\begin{lstlisting}
# Core packages
install.packages("tidyverse")   # Data manipulation and visualization
install.packages("here")        # File path management

# Summary statistics
install.packages("psych")       # describe() function
install.packages("Hmisc")       # More descriptive stats

# Visualization
install.packages("ggplot2")     # Already in tidyverse, but can load separately
install.packages("gridExtra")   # Arrange multiple plots

# Statistical tests
# Base R includes t.test(), wilcox.test(), aov(), chisq.test(), cor.test()

# Assumption checking
install.packages("car")         # Levene's test, Q-Q plots

# RMarkdown (if using)
install.packages("rmarkdown")
install.packages("knitr")
\end{lstlisting}

\section*{Getting Help}

If your group encounters difficulties:

\subsection*{Office Hours}
\begin{itemize}
\item Instructor office hours: [Days/times listed on syllabus]
\item TA office hours: [Days/times listed on syllabus]
\item Special project office hours Week 7.5 for dataset approval/questions
\end{itemize}

\subsection*{Canvas Discussion Board}
\begin{itemize}
\item Post general questions (not dataset-specific answers)
\item Help other groups when you can
\item Instructor monitors regularly
\end{itemize}

\subsection*{Email}
\begin{itemize}
\item For private group concerns
\item For dataset approval requests
\item For deadline extensions (must request BEFORE due date)
\end{itemize}

\subsection*{Resources}
\begin{itemize}
\item Course textbook: Pagano \& Gauvreau, \textit{Principles of Biostatistics}
\item R documentation: \url{https://www.r-project.org/}
\item RMarkdown guide: \url{https://rmarkdown.rstudio.com/}
\item Course Canvas page: lecture notes, example code
\end{itemize}

\section*{Academic Integrity}

\begin{tcolorbox}[colback=gray!10,colframe=black,boxrule=1pt]
\textbf{Collaboration Policy:}

\begin{itemize}
\item \textbf{Within your group}: Full collaboration is expected and encouraged
\item \textbf{Between groups}: You may discuss general concepts, but not share code or specific analyses
\item \textbf{Outside resources}: You may consult textbooks, online tutorials, AI tools, but must cite sources
\item \textbf{Plagiarism}: Copying from other groups, previous years, or uncited sources is prohibited
\end{itemize}

\textbf{What this means:}
\begin{itemize}
\item It's OK to say: ``Our group is using ANOVA for our main analysis. What are you using?''
\item It's NOT OK to say: ``Can I see your R code for the ANOVA?''
\item It's OK to use AI to debug code or explain concepts
\item It's NOT OK to have AI write your entire analysis
\end{itemize}

All group members are responsible for ensuring academic integrity. Violations will result in zero points for the assignment and possible further consequences per university policy.
\end{tcolorbox}

\vspace{12pt}

\begin{center}
\textbf{\Large Good luck with your project!}\\
\vspace{6pt}
\textit{This is an opportunity to apply your biostatistics knowledge to real-world data.}\\
\textit{Ask questions, collaborate with your team, and enjoy the process of discovery.}
\end{center}

\end{document}
